%% kapa: exam=1 round=19 year=2008 period=2 form=A subject=economics subject_ko=경제원론 seq=15 number=15
%% source: Q-Net 감정평가사 기출문제 (cst00309) artlSeq=1001962 file=extracted/1008/1001962/2008년 제19회 감정평가사 1차시험 기출문제/2교시.hwp
%% conversion: kordoc 4.24.51 -> tools/gen1.py (auto); review=proofread-sonnet; flags=none
\begin{question}{15}
정부 정책과 관련된 설명 중 옳은 것을 모두 고르면?
\begin{framedtext}
ㄱ. 시간 불일치(time inconsistency)의 문제를 완화하려면 준칙에 따른 정책을 수행해야 한다. \par
ㄴ. 균형을 회복하는 시장의 조정과정이 빠를수록 재량적 정책이 선호된다. \par
ㄷ. 정책 시차가 길고 정책의 효과가 불확실할수록 재량적 정책이 선호된다. \par
\end{framedtext}
\begin{choices}
  \item ㄱ
  \item ㄴ
  \item ㄷ
  \item ㄱ, ㄴ
  \item ㄱ, ㄷ
\end{choices}
\end{question}
\answer{1}
